\section{Appendix: Causal interpretation and benchmark proofs}

The causal interpretation and the benchmark comparisons concern the same observable class in \cref{obj:def:model}, defined by the public overlap restriction. The potential-outcome statement describes compatible extensions of each law in that class. The benchmark comparisons concern the information supplied to estimation rules under the sampling experiment in \cref{obj:ass:data-product}.

\subsection{Compatible extensions and identification}

The existence statement uses a conditional construction with independent binary treatment and potential outcomes. We first specify the probability law used in that construction.

\begin{definitionv}{synth_7}[Bernoulli law]
For \(p\in[0,1]\), \(\operatorname{Bernoulli}(p)\) is the probability law on \(\{0,1\}\) assigning probability \(p\) to \(1\) and probability \(1-p\) to \(0\).
\end{definitionv}

This law expresses binary conditional distributions directly in terms of their success probabilities. In the following construction, the observed propensity and outcome regressions determine those probabilities within each occupied covariate cell.

\begin{lemmav}{lem:causal-realization}[Causal realization and identification]
Let \(d\) be a natural number, \(\epsilon\) a real number, and \(P\) a probability law of the complete observation \((X,A,Y)\). Suppose:
\begin{itemize}
\item \textbf{(Covariate alphabet.)} \(d\ge 2\).
\item \textbf{(Overlap range.)} \(0<\epsilon\le 1/4\).
\item \textbf{(Model membership.)} \(P\in\mathcal M_{d,\epsilon}\), as defined in \cref{obj:def:model}.
\end{itemize}
Then \(P\) admits a joint probability law \(H\) of \((X,A,Y,Y_0,Y_1)\), with binary potential outcomes \(Y_0,Y_1\), whose \((X,A,Y)\) marginal equals \(P\) and which satisfies \cref{obj:ass:consistency,obj:ass:exchangeability}. Call precisely these extensions compatible.

One compatible extension is given by
\[
\Pr_H(X=j)=p_j(P),
\]
and, for every occupied cell \(j\), by the conditional law
\[
\mathcal L_H(A,Y_0,Y_1\mid X=j)
=
\operatorname{Bernoulli}\!\bigl(e_j(P)\bigr)
\otimes
\operatorname{Bernoulli}\!\bigl(\mu_{0j}(P)\bigr)
\otimes
\operatorname{Bernoulli}\!\bigl(\mu_{1j}(P)\bigr),
\qquad Y=Y_A.
\]
For every compatible extension \(H\), the causal contrast equals the target defined in \cref{obj:def:target}:
\[
\mathbb E_H[Y_1-Y_0]=\tau(P).
\]
\end{lemmav}

The product construction supplies an existence witness for every observable law in the stated class. It chooses conditional independence between the two potential outcomes as well as between treatment and the potential-outcome pair. Compatibility itself requires the observed marginal, consistency, and conditional treatment exchangeability; the identifying equality applies to every extension satisfying those requirements. Thus the causal contrast is invariant across compatible choices of the joint potential-outcome distribution.

The public overlap restriction gives both treatment arms positive probability in every occupied covariate cell. Under a compatible extension, consistency and conditional exchangeability identify each conditional potential-outcome mean with its corresponding observed outcome regression. The target in \cref{obj:def:target} aggregates their contrast using the population covariate weights. Null covariate cells have zero weight under both the observable law and its extensions. These roles keep the observable restriction in \cref{obj:ass:overlap} distinct from the identification conditions imposed on the extension in \cref{obj:ass:consistency,obj:ass:exchangeability}.

\subsection{Two-cell and exact-marginal comparisons}

The benchmark functions \(b(n,\epsilon)\) and \(f(n,m,d)\) are defined in \cref{obj:prop:benchmark-recovery}: the former is the capped inverse rare-arm outcome scale, and the latter combines inverse complete-sample precision with a many-cell term at fixed overlap. The numerical rate \(r(n,m,d,\epsilon)\) is defined in \cref{obj:def:rate-handle}. These definitions distinguish the numerical comparisons from the minimax risks in \cref{obj:def:risk,obj:def:known-risk}.

For a two-cell covariate alphabet, \cref{obj:prop:benchmark-recovery} gives rare-label order uniformly over the auxiliary budget and the allowed overlap range. The uniform risk comparison in \cref{obj:thm:uniform-annotation-frontier} supplies the connection between the numerical rate and attainable precision. The matching lower bound in \cref{obj:thm:rare-label-floor} locates the persistent difficulty in the treated outcomes of the complete records.

The exact-marginal comparison has the same rare-label order for every allowed alphabet size. In this experiment, the population covariate weights and treatment propensities are supplied through the marginal table, as specified in \cref{obj:def:known-risk}. The common-marginal alternatives in \cref{obj:thm:rare-label-floor} retain their target separation when this table is supplied. Together, the two comparisons in \cref{obj:prop:benchmark-recovery} characterize experiments whose minimax precision is governed by complete-record outcome information, with numerical comparison constants uniform over their stated public indices.

At each fixed positive overlap floor, the comparison with \(f(n,m,d)\) in \cref{obj:prop:benchmark-recovery} expresses the rate in the conventional complete-sample and total-sample scales. Its constants may depend on that fixed floor. The two-cell and exact-marginal comparisons retain uniform constants as the floor varies within the stated range.

\subsection{The fixed-parameter experiment limit}

The exact-marginal risk provides the limiting benchmark for increasing auxiliary information. For \(n\ge1\), \(m\ge1\), \(d\ge2\), and \(0<\epsilon\le1/4\), \cref{obj:prop:benchmark-recovery} bounds the nonnegative difference between the two-channel and exact-marginal risks by
\[
\min\left\{1,6(n+2)\sqrt{\frac{2d}{m}}\right\}.
\]
This comparison concerns the minimax risk itself. At fixed \(n,d,\epsilon\), its dependence on the auxiliary sample size quantifies convergence to the exact-marginal experiment.

The same proposition establishes that bounded rare-arm outcome information gives a uniformly positive risk floor over all auxiliary budgets and allowed alphabet sizes. The experiment limit and this floor therefore describe complementary aspects of auxiliary information: increasing auxiliary records approaches the precision available with the population marginal supplied, while the rare-label benchmark characterizes the outcome information carried by the fixed complete sample.
